\subsection{\label{sec:psu3}PSU3}

\subsubsection{Action}

The link variables are represented by matrices $U_\mu(n)\in SU(3)$, but the
action is chosen to be insensitive to the center of $SU(3)$.  For
$z\in Z_3$ the adjoint representation satisfies
\begin{equation}
D_{\rm adj}(zU)=D_{\rm adj}(U).
\end{equation}
Equivalently, the character identity following from
$\mathbf{3}\otimes\overline{\mathbf{3}}=\mathbf{1}\oplus\mathbf{8}$ is
\begin{equation}
\operatorname{Tr}_{\rm adj} U
=|\operatorname{Tr}_{F}U|^2-1 .
\end{equation}
The adjoint Wilson plaquette action can therefore be evaluated using only
fundamental $3\times3$ matrices:
\begin{equation}
S_A=-\beta_A\sum_n\sum_{\mu<\nu}\frac{1}{8}
\left(|\operatorname{Tr}U_{\mu\nu}(n)|^2-1\right),
\end{equation}
where $1/8$ is the normalization by the dimension of the adjoint
representation and
\begin{equation}
U_{\mu\nu}(n)=U_\mu(n)U_\nu(n+\hat\mu)
U_\mu^\dagger(n+\hat\nu)U_\nu^\dagger(n).
\end{equation}

This remains a one-plaquette action.  The squared quantity is the modulus of
the fundamental trace; it is neither $U_{\mu\nu}^2$ nor a geometrically
doubled plaquette.  It appears only because the adjoint character is being
written in terms of fundamental matrices.

Indeed, the links may be transformed independently as
\begin{equation}
U_\mu(n)\longrightarrow z_\mu(n)U_\mu(n),
\qquad z_\mu(n)\in Z_3.
\end{equation}
For a plaquette this gives
\begin{equation}
U_{\mu\nu}(n)\longrightarrow z_{\mu\nu}(n)U_{\mu\nu}(n),
\qquad z_{\mu\nu}(n)\in Z_3,
\end{equation}
and hence
\begin{equation}
|\operatorname{Tr}U_{\mu\nu}(n)|^2
\longrightarrow
|z_{\mu\nu}(n)|^2|\operatorname{Tr}U_{\mu\nu}(n)|^2
=|\operatorname{Tr}U_{\mu\nu}(n)|^2.
\end{equation}
Thus center-related $SU(3)$ link matrices represent the same physical
$PSU(3)=SU(3)/Z_3$ link variable.  The continuous Lie algebra is nevertheless
the usual $\mathfrak{su}(3)$ algebra, so the HMC momenta and tangent-space
projection are the same as for an $SU(3)$ link model.

The additive constant in $S_A$ does not affect molecular dynamics.  Defining
\begin{equation}
\widetilde\beta=\frac{\beta_A}{8},
\end{equation}
the dynamic part used below is
\begin{equation}
\boxed{
S_{\rm dyn}=-\widetilde\beta\sum_n\sum_{\mu<\nu}
|\operatorname{Tr}U_{\mu\nu}(n)|^2 .
}
\end{equation}

\subsubsection{HMC Force}

Fix one link and abbreviate
\begin{equation}
U=U_\mu(n).
\end{equation}
For every $\nu\ne\mu$, two plaquettes in the $\mu\nu$ plane contain this
link.  After cyclically rotating the trace, both can be written with $U$ as
their first factor:
\begin{align}
W_{+}(\mu,\nu,n)&=U B_{+}(\mu,\nu,n),\\
W_{-}(\mu,\nu,n)&=U B_{-}(\mu,\nu,n),
\end{align}
where
\begin{align}
B_{+}(\mu,\nu,n)
&=U_\nu(n+\hat\mu)U_\mu^\dagger(n+\hat\nu)U_\nu^\dagger(n),\\
B_{-}(\mu,\nu,n)
&=U_\nu^\dagger(n+\hat\mu-\hat\nu)
  U_\mu^\dagger(n-\hat\nu)U_\nu(n-\hat\nu).
\end{align}
Their traces satisfy
\begin{align}
t_{+}&\equiv\operatorname{Tr}W_{+}
      =\operatorname{Tr}U_{\mu\nu}(n),\\
t_{-}&\equiv\operatorname{Tr}W_{-}
      =\operatorname{Tr}U_{\nu\mu}(n-\hat\nu).
\end{align}

We first calculate the raw component derivative by regarding $U$ and
$U^\dagger$ as independent variables.  In particular, $U^\dagger$ is held
fixed in $\partial S/\partial U$.  For a generic path $W=UB$,
\begin{equation}
\operatorname{Tr}W=U^l{}_k B^k{}_l,
\end{equation}
and therefore
\begin{equation}
\frac{\partial\operatorname{Tr}W}{\partial U^j{}_i}
=\delta^l{}_j\delta^i{}_k B^k{}_l
=B^i{}_j.
\end{equation}
Since the conjugate trace is held fixed in this raw derivative,
\begin{equation}
\left.
\frac{\partial |\operatorname{Tr}W|^2}{\partial U^j{}_i}
\right|_{U^\dagger}
=(\operatorname{Tr}W)^*B^i{}_j.
\end{equation}
It follows that the raw derivative of the action with respect to the chosen
link is
\begin{equation}
\boxed{
\left.
\left(\frac{\partial S}{\partial U_\mu(n)}\right)^i{}_j
\right|_{U^\dagger}
=-\widetilde\beta\sum_{\nu\ne\mu}
\left[t_{+}^* B_{+}{}^i{}_j+t_{-}^* B_{-}{}^i{}_j\right].
}
\end{equation}
Equivalently, in terms of the original plaquette notation,
\begin{equation}
\boxed{
\left.
\left(\frac{\partial S}{\partial U_\mu(n)}\right)^i{}_j
\right|_{U^\dagger}
=-\widetilde\beta\sum_{\nu\ne\mu}
\left[
\operatorname{Tr}U_{\nu\mu}(n)B_{+}{}^i{}_j
+\operatorname{Tr}U_{\mu\nu}(n-\hat\nu)B_{-}{}^i{}_j
\right].
}
\end{equation}

The coefficient is $-\widetilde\beta$, rather than
$-\widetilde\beta/2$.  One may write the action as
\begin{equation}
S_{\rm dyn}=-\frac{\widetilde\beta}{2}
\sum_n\sum_{\mu\ne\nu}
\operatorname{Tr}U_{\mu\nu}(n)\operatorname{Tr}U_{\nu\mu}(n),
\end{equation}
but the ordered sum contains both $(\mu,\nu)$ and $(\nu,\mu)$.  Their
identical derivatives cancel the explicit factor $1/2$.

To compare the component derivative with the force field stored by the
action, define the code staple by
\begin{equation}
\Sigma_P=B_P^\dagger,
\qquad U_P=U\Sigma_P^\dagger=UB_P.
\end{equation}
The raw field is
\begin{equation}
Y_\mu(n)=-\widetilde\beta\sum_{P\ni(n,\mu)}
\operatorname{Tr}(U_P)\Sigma_P(n).
\end{equation}
Consequently, the tensor-element result is precisely
\begin{equation}
\boxed{
\left.\frac{\partial S}{\partial U_\mu(n)}\right|_{U^\dagger}
=Y_\mu^\dagger(n),
}
\end{equation}
with the index placement understood as in the component equations above.

For completeness, the same result may be expressed directly as an
$\mathfrak{su}(3)$ variation.  Let $A$ be traceless and anti-Hermitian and
vary the link from the left,
\begin{equation}
U\longrightarrow e^A U,
\qquad \delta U=AU.
\end{equation}
For one oriented plaquette $U_P=UB_P$ this gives
\begin{equation}
\delta S_P=-\widetilde\beta\operatorname{Tr}\left\{
A\left[
\operatorname{Tr}(U_P^\dagger)U_P
-\operatorname{Tr}(U_P)U_P^\dagger
\right]\right\}.
\end{equation}
Defining the mathematical force by
$\delta S=-\operatorname{Tr}(AF_\mu(n))$, one obtains
\begin{equation}
\boxed{
F_\mu(n)=\widetilde\beta\sum_{P\ni(n,\mu)}X_P,
\qquad
X_P=\operatorname{Tr}(U_P^\dagger)U_P
-\operatorname{Tr}(U_P)U_P^\dagger .
}
\end{equation}
The matrix $X_P$ is automatically traceless and anti-Hermitian.  This last
form is only the tangent-space version of the tensor derivative above; the
two expressions are not independent force formulas.

\subsubsection{CLGLib force convention}

The matrix $F_\mu(n)$ defined above is the mathematical force in
\begin{equation}
\delta S=-\operatorname{Tr}\bigl(A F_\mu(n)\bigr).
\end{equation}
CLGLib stores a different, but directly related, matrix in its force field.
The gauge momentum is anti-Hermitian and has kinetic energy
\begin{equation}
K=-\operatorname{Tr}(P^2).
\end{equation}
With the molecular-dynamics equation $\dot U=PU$, one has
\begin{equation}
\dot S=-\operatorname{Tr}(P F),
\qquad
\dot K=-2\operatorname{Tr}(P\dot P).
\end{equation}
Energy conservation therefore requires
\begin{equation}
\boxed{
\dot P=-\frac{1}{2}F.
}
\end{equation}
Thus the matrix added to the momentum by the integrator is
\begin{equation}
F^{\rm code}_\mu(n)=-\frac{1}{2}F_\mu(n).
\end{equation}

For any matrix $M$, CLGLib's operation \texttt{TA()} is
\begin{equation}
{\rm TA}(M)=\frac{1}{2}(M-M^\dagger)
-\frac{1}{3}\operatorname{Tr}\left[\frac{1}{2}(M-M^\dagger)\right]\mathbf 1.
\end{equation}
The action kernel stores the raw matrix $Y_\mu(n)$, after which the common
integrator performs
\begin{equation}
\boxed{
F^{\rm code}_\mu(n)
={\rm TA}\Bigl(U_\mu(n)Y_\mu^\dagger(n)\Bigr).
}
\end{equation}
There is no additional factor $-1/2$ outside \texttt{TA()}; the factor
$1/2$ is already part of the anti-Hermitian projection.

Indeed, using
\begin{equation}
Y_\mu(n)=-\widetilde\beta\sum_{P\ni(n,\mu)}
\operatorname{Tr}(U_P)\Sigma_P(n),
\qquad
U_P=U_\mu(n)\Sigma_P^\dagger(n),
\end{equation}
one obtains
\begin{align}
U_\mu(n)Y_\mu^\dagger(n)
&=-\widetilde\beta\sum_P
  \operatorname{Tr}(U_P^\dagger)
  U_\mu(n)\Sigma_P^\dagger(n)\nonumber\\
&=-\widetilde\beta\sum_P
  \operatorname{Tr}(U_P^\dagger)U_P.
\end{align}
Taking the traceless anti-Hermitian part gives
\begin{equation}
\boxed{
F^{\rm code}_\mu(n)
=-\frac{\widetilde\beta}{2}\sum_{P\ni(n,\mu)}X_P
=-\frac{1}{2}F_\mu(n).
}
\end{equation}

\subsubsection{CUDA implementation}

The implementation in \texttt{CActionGaugePlaquettePSU3.cu} follows the
preceding equations in two stages.

For each link, \texttt{\_kernelForce\_PSU3} loops over the surrounding
three-link paths supplied by the staple cache.  For each path it performs
\begin{align}
\Sigma_P&\leftarrow\text{cached three-link path},\\
U_P&\leftarrow U_\mu(n)\Sigma_P^\dagger,\\
t_P&\leftarrow\operatorname{Tr}(U_P),\\
Y_\mu(n)&\leftarrow
Y_\mu(n)-\widetilde\beta\,t_P\Sigma_P.
\end{align}
After all actions have contributed their raw fields, the common integrator
performs \texttt{LeftMul(gauge, FALSE, TRUE)}, which replaces the raw field by
$U_\mu Y_\mu^\dagger$, and then calls \texttt{TA()}.

The trace $t_P$ must be computed separately for every plaquette.  Although
the trace itself is linear, the force contains the weighted sum
\begin{equation}
\sum_P \operatorname{Tr}(U_P)\Sigma_P,
\end{equation}
whose weight depends on $P$.  It cannot be reconstructed from one trace of a
pre-summed staple.

For the $SU(N)$ fields supported by this class, the normalization uses the
adjoint dimension:
\begin{equation}
\texttt{Beta}=\beta_A,
\qquad
\widetilde\beta=\frac{\beta_A}{N^2-1}
=\begin{cases}
\beta_A/8, & SU(3),\\[2pt]
\beta_A/3, & SU(2).
\end{cases}
\end{equation}

There are two implementation qualifications:
\begin{itemize}
\item Raw gauge forces in CLGLib are additive because several actions may
contribute to the same link.  The kernel therefore \emph{adds} its raw
contribution to the shared force field via \texttt{\_add(pForceData[uiLinkIndex], res)},
exactly as the standard plaquette kernel does; the field is zeroed once by the
integrator before any action is evaluated.
\item The coupling is normalized by the adjoint dimension $N^2-1$, so the same
class works for $SU(3)$ ($\widetilde\beta=\beta_A/8$) and $SU(2)$
($\widetilde\beta=\beta_A/3$).  Other matrix sizes are reported as unsupported
at initialization and are rejected by the force and energy routines.
\end{itemize}

\subsubsection{Coupling to fermions}

At the level of the HMC program, this gauge action can be combined with a
fermion action using the same $SU(3)$ link matrices: each action contributes
to the total energy and adds its force to the shared force field.  However,
the resulting theory is a genuine $PSU(3)$ gauge theory only if the fermions
are neutral under the center $Z_3$.

For example, a fundamental-fermion hopping term contains
\begin{equation}
\overline\psi(n)U_\mu(n)\psi(n+\hat\mu).
\end{equation}
Under the independent link-center transformation
$U_\mu(n)\to z_\mu(n)U_\mu(n)$, this term acquires the factor $z_\mu(n)$ and
is not invariant.  A fundamental $\mathbf 3$ is therefore a representation of
$SU(3)$, but not a representation of $PSU(3)=SU(3)/Z_3$.

Consequently, combining the adjoint plaquette action with ordinary
fundamental $SU(3)$ fermions is a valid lattice model, but it should be
interpreted as an $SU(3)$ gauge theory with an adjoint gauge action.  The
fermion action removes the independent local $Z_3$ identification that made
the pure-gauge sector center blind.  To preserve the $PSU(3)$ interpretation,
one should use center-neutral, zero-triality matter, such as adjoint fermions.
Introducing fundamental matter while retaining a quotient-gauge
interpretation requires additional structure and is not achieved by simply
adding a standard fundamental-fermion action.
